% GENERATED FILE. Edit canonical lessons, then run npm run generate:books.
% canonical-source-hash: fnv1a64:359aaf3867723711
% canonical-lessons: HI-C442-sikhar, HI-C442-tila, HI-C442-dher, HI-C442-tinka, HI-C442-lurhkana

\chapter{Peaks and Mounds}
\label{ch:hi-peaks-and-mounds}
\hlchaptermodality{442}

\begin{chapteropening}
\textbf{By the end of this chapter:} \emph{I can talk about a summit, a dune, a heap and a straw, and stones rolled down a hill.}

Complete the last lesson of chapter 442: I can talk about a summit, a dune, a heap and a straw, and stones rolled down a hill.
\end{chapteropening}

\section[\texorpdfstring{\hi{शिखर} (śikhar)}{शिखर (śikhar)}]{\hi{शिखर} (śikhar) — a peak, a summit}

\label{lesson:HI-C442-sikhar}



\begin{quote}
Before the new one: say the Hindi for the horizon, then the Hindi for dusk.
\end{quote}

\subsection*{You'll want to know: \hi{शिखर}}
\textbf{\hi{शिखर}} — \emph{śikhar} — \textquotedblleft{}a peak, a summit\textquotedblright{}.

\begin{quote}
\hi{हम} \hi{सुबह} \hi{तक} \hi{पहाड़} \hi{के} \hi{शिखर} \hi{पर} \hi{पहुँचे।}
\end{quote}

\begin{quote}
\emph{ham subah tak pahāṛ ke śikhar par pahũce.}
\end{quote}

\begin{quote}
We reached the top of the mountain by morning.
\end{quote}

Masculine.

\begin{cousinweb}
\textbf{\hi{शिखर}} (\emph{śikhar}) is a tatsama, Sanskrit \emph{śikhara}, \textquotedblleft{}pointed, crested\textquotedblright{} and so a peak (Monier-Williams); the spire of a temple is its \hi{शिखर} too.
\end{cousinweb}

\begin{grammarlens}[title={What you've built}]
One more word for getting around.
\end{grammarlens}

\subsection*{Your turn}
\begin{itemize}
\raggedright
  \item \emph{Say it:} \emph{śikhar}
  \item \emph{Say it:} \emph{śikhar}, once more
  \item \emph{Say it:} \emph{ham subah tak pahāṛ ke śikhar par pahũce.} — We reached the top of the mountain by morning.
  \item \emph{Recall:} say the Hindi for the horizon, then the Hindi for dusk, then say \emph{śikhar} again
\end{itemize}

\begin{tcolorbox}[breakable,colback=teal!4,colframe=teal!35!black,title={Before you move on}]
What does \hi{शिखर} mean? (\textquotedblleft{}A peak, a summit\textquotedblright{}.) Say it once more.
\end{tcolorbox}
\section[\texorpdfstring{\hi{टीला} (ṭīlā)}{टीला (ṭīlā)}]{\hi{टीला} (ṭīlā) — a mound; a dune}

\label{lesson:HI-C442-tila}



\begin{quote}
Before the new one: say the Hindi for dusk, then the Hindi for a summit.
\end{quote}

\subsection*{You'll want to know: \hi{टीला}}
\textbf{\hi{टीला}} — \emph{ṭīlā} — \textquotedblleft{}a mound; a dune\textquotedblright{}.

\begin{quote}
\hi{रेगिस्तान} \hi{में} \hi{रेत} \hi{के} \hi{ऊँचे} \hi{टीले} \hi{हैं।}
\end{quote}

\begin{quote}
\emph{registān meṁ ret ke ū̃ce ṭīle haiṁ.}
\end{quote}

\begin{quote}
There are high sand dunes in the desert.
\end{quote}

Masculine.

\begin{cousinweb}
\textbf{\hi{टीला}} (\emph{ṭīlā}) goes back to \emph{*ṭalla}, \emph{*ṭilla}, \textquotedblleft{}a heap\textquotedblright{}, forms CDIAL reconstructs (CDIAL 5451), which lists Hindi \emph{ṭīlā}, \textquotedblleft{}a hillock, a mound\textquotedblright{}.
\end{cousinweb}

\begin{grammarlens}[title={What you've built}]
One more word for getting around.
\end{grammarlens}

\subsection*{Your turn}
\begin{itemize}
\raggedright
  \item \emph{Say it:} \emph{ṭīlā}
  \item \emph{Say it:} \emph{ṭīlā}, once more
  \item \emph{Say it:} \emph{registān meṁ ret ke ū̃ce ṭīle haiṁ.} — There are high sand dunes in the desert.
  \item \emph{Recall:} say the Hindi for dusk, then the Hindi for a summit, then say \emph{ṭīlā} again
\end{itemize}

\begin{tcolorbox}[breakable,colback=teal!4,colframe=teal!35!black,title={Before you move on}]
What does \hi{टीला} mean? (\textquotedblleft{}A mound; a dune\textquotedblright{}.) Say it once more.
\end{tcolorbox}
\section[\texorpdfstring{\hi{ढेर} (ḍher)}{ढेर (ḍher)}]{\hi{ढेर} (ḍher) — a heap, a pile}

\label{lesson:HI-C442-dher}



\begin{quote}
Before the new one: say the Hindi for a summit, then the Hindi for a mound.
\end{quote}

\subsection*{You'll want to know: \hi{ढेर}}
\textbf{\hi{ढेर}} — \emph{ḍher} — \textquotedblleft{}a heap, a pile\textquotedblright{}.

\begin{quote}
\hi{गली} \hi{के} \hi{कोने} \hi{पर} \hi{कूड़े} \hi{का} \hi{ढेर} \hi{है।}
\end{quote}

\begin{quote}
\emph{galī ke kone par kūṛe kā ḍher hai.}
\end{quote}

\begin{quote}
There is a heap of rubbish at the corner of the lane.
\end{quote}

Masculine.

\begin{cousinweb}
\textbf{\hi{ढेर}} (\emph{ḍher}) goes back to \emph{*ḍhēra}, \textquotedblleft{}a lump, a heap\textquotedblright{}, a form CDIAL reconstructs (CDIAL 5599), which lists Hindi \emph{ḍher}, \textquotedblleft{}heap\textquotedblright{}.
\end{cousinweb}

\begin{grammarlens}[title={What you've built}]
One more word for getting around.
\end{grammarlens}

\subsection*{Your turn}
\begin{itemize}
\raggedright
  \item \emph{Say it:} \emph{ḍher}
  \item \emph{Say it:} \emph{ḍher}, once more
  \item \emph{Say it:} \emph{galī ke kone par kūṛe kā ḍher hai.} — There is a heap of rubbish at the corner of the lane.
  \item \emph{Recall:} say the Hindi for a summit, then the Hindi for a mound, then say \emph{ḍher} again
\end{itemize}

\begin{tcolorbox}[breakable,colback=teal!4,colframe=teal!35!black,title={Before you move on}]
What does \hi{ढेर} mean? (\textquotedblleft{}A heap, a pile\textquotedblright{}.) Say it once more.
\end{tcolorbox}
\section[\texorpdfstring{\hi{तिनका} (tinkā)}{तिनका (tinkā)}]{\hi{तिनका} (tinkā) — a straw; a blade of grass}

\label{lesson:HI-C442-tinka}



\begin{quote}
Before the new one: say the Hindi for a mound, then the Hindi for a heap.
\end{quote}

\subsection*{You'll want to know: \hi{तिनका}}
\textbf{\hi{तिनका}} — \emph{tinkā} — \textquotedblleft{}a straw; a blade of grass\textquotedblright{}.

\begin{quote}
\hi{हवा} \hi{में} \hi{तिनके} \hi{उड़} \hi{रहे} \hi{थे।}
\end{quote}

\begin{quote}
\emph{havā meṁ tinke uṛ rahe the.}
\end{quote}

\begin{quote}
Bits of straw were flying in the wind.
\end{quote}

Masculine.

\begin{cousinweb}
\textbf{\hi{तिनका}} (\emph{tinkā}) is a tadbhava of Sanskrit \emph{tṛṇa}, \textquotedblleft{}grass, a blade of grass, a straw\textquotedblright{} (CDIAL 5906), which lists Hindi \emph{tinkā}: in sayings it is the smallest thing there is.
\end{cousinweb}

\begin{grammarlens}[title={What you've built}]
One more word for getting around.
\end{grammarlens}

\subsection*{Your turn}
\begin{itemize}
\raggedright
  \item \emph{Say it:} \emph{tinkā}
  \item \emph{Say it:} \emph{tinkā}, once more
  \item \emph{Say it:} \emph{havā meṁ tinke uṛ rahe the.} — Bits of straw were flying in the wind.
  \item \emph{Recall:} say the Hindi for a mound, then the Hindi for a heap, then say \emph{tinkā} again
\end{itemize}

\begin{tcolorbox}[breakable,colback=teal!4,colframe=teal!35!black,title={Before you move on}]
What does \hi{तिनका} mean? (\textquotedblleft{}A straw; a blade of grass\textquotedblright{}.) Say it once more.
\end{tcolorbox}
\section[\texorpdfstring{\hi{लुढ़काना} (luṛhkānā)}{लुढ़काना (luṛhkānā)}]{\hi{लुढ़काना} (luṛhkānā) — to roll (something)}

\label{lesson:HI-C442-lurhkana}



\begin{quote}
Before the new one: say the Hindi for a heap, then the Hindi for a straw.
\end{quote}

\subsection*{You'll want to know: \hi{लुढ़काना}}
\textbf{\hi{लुढ़काना}} — \emph{luṛhkānā} — \textquotedblleft{}to roll (something)\textquotedblright{}.

\begin{quote}
\hi{बच्चों} \hi{ने} \hi{टीले} \hi{से} \hi{पत्थर} \hi{नीचे} \hi{लुढ़काए।}
\end{quote}

\begin{quote}
\emph{baccoṁ ne ṭīle se patthar nīce luṛhkāe.}
\end{quote}

\begin{quote}
The children rolled stones down the mound.
\end{quote}

Past, with \hi{ने} (he or she …): \hi{उसने} \hi{ढोल} \hi{लुढ़काया}, \emph{usne ḍhol luṛhkāyā}.

\begin{cousinweb}
\textbf{\hi{लुढ़काना}} (\emph{luṛhkānā}) is the causative partner of \hi{लुढ़कना}, \textquotedblleft{}to roll\textquotedblright{}, from Sanskrit \emph{luṭhati}, \textquotedblleft{}rolls\textquotedblright{} (CDIAL 11079): \hi{लुढ़कना} is to roll, \hi{लुढ़काना} to roll a thing.
\end{cousinweb}

\begin{grammarlens}[title={What you've built}]
One more word for getting around.
\end{grammarlens}

\subsection*{Your turn}
\begin{itemize}
\raggedright
  \item \emph{Say it:} \emph{luṛhkānā}
  \item \emph{Say it:} \emph{luṛhkānā}, once more
  \item \emph{Say it:} \emph{baccoṁ ne ṭīle se patthar nīce luṛhkāe.} — The children rolled stones down the mound.
  \item \emph{Recall:} say the Hindi for a heap, then the Hindi for a straw, then say \emph{luṛhkānā} again
\end{itemize}

\begin{tcolorbox}[breakable,colback=teal!4,colframe=teal!35!black,title={Before you move on}]
What does \hi{लुढ़काना} mean? (\textquotedblleft{}To roll (something)\textquotedblright{}.) Say it once more.
\end{tcolorbox}
